\section{Расчётная задача 8}

Пусть $V = \RR^3$ --- векторное пространство над~$\RR$, $q \in S^2V$ --- квадратичная форма, $\ee = (e_1, e_2, e_3) \subset V$ --- ортонормированный базис.

Требуется:
\begin{enumerate}[1)]
	\item привести~$q$ к каноническому виду методом Лагранжа; выписать соответсвующее преобразование координат; а также выписать матрицы~$q$ в исходном и каноническом базисах; проверить, является ли преобразование координат ортогональным;
	\item проверить с помощью критерия Сильвестра, является ли  квадратичная форма знакоопределённой;
	\item классифицировать~$q$ с использованием индексов инерции.
\end{enumerate}

Форма~$q$ в базисе~$\ee$ имеет вид
\[
	q =
	e_1 \vee e_1
	-
	5
	e_2 \vee e_2
	+
	2e_1 \vee e_2
	-
	4e_3 \vee e_4
	+
	2e_1 \vee e_4
	.
\]
Если~$x = \ee (x_1, x_2, x_3, x_4)$, то
\[
	q(x)
	=
	(x_1)^2 - 5(x_2)^2 + 2x_1 x_2 - 4 x_3 x_4 + 2x_1 x_4
	.
\]
(В этой задаче мы будем писать индексы у координат вектора \emph{снизу}. Это не слишком правильно, но типографически проще.)

\begin{proofm}
	Заметим, что~$q$ в базисе~$\ee$ имеет матрицу
	\[
		\Qe =
		\begin{pmatrix}
			1 & 1 & 0 & 1 \\
			1 & -5 & 0 & 0 \\
			0 & 0 & 0 & -2 \\
			1 & 0 & -2 & 0
		\end{pmatrix}
		.
	\]
	Это значит, что если
	\[
		\crd{x}_{\ee} =
		\Vect{x_1 \\ x_2 \\ x_3 \\ x_4}
	\]
	координаты вектора~$x$ в~базисе~$\ee$, то
	\[
		q(x) =
		\crd{x}_{\ee}^\top
		\Qe
		\crd{x}_{\ee}
	\]
	в смысле умножения матриц.

	Далее, если $\ss = \ee S$ --- замена координат, то
	\[
		x =
		\ee \crd{x}_{\ee} =
		\ss \crd{x}_{\ss} =
		\ee S \crd{x}_{\ss}
		,
	\]
	откуда
	\[
		\crd{x}_{\ee} = S \crd{x}_{\ss},
	\]
	и
	\[
		q(x) =
		\bigl(
			S \crd{x}_{\ss}
		\bigr)^\top
		\Qe
		\bigl(
			S \crd{x}_{\ss}
		\bigr)
		=
		\crd{x}_{\ss}^\top
		\cdot
		\bigl(
			\underbrace{S^\top \Qe S}_{\Qs}
		\bigr)
		\cdot
		\crd{x}_{\ss}
		,
	\]
	поэтому
	\[
		\Qs = S^{\top} \Qe S,
		\qquad
		\text{где $\ss = \ee S$.}
	\]
	
	Итак, наша квадратичная форма~$q$ в базисе~$\ee$ в координатах выглядит так:
	\[
		q(x)
		=
		x_1^2 - 5x_2^2 + 2x_1 x_2 - 4 x_3 x_4 + 2x_1 x_4
		.
	\]

	Выделим полный квадрат для переменной, содержащей~$x_1$:
	\begin{align*}
		q(x)
		&=
		x_1^2
		\noatt{ {} - 5x_2^2 }
		+
		2x_1 x_2 
		\noatt{ {} - 4 x_3 x_4 }
		+ 2x_1 x_4
		\nlinea{=}
		\bigl(
			x_1^2  + 2(x_2 + x_4)x_1 
		\bigr)
		\noatt{ {}
			- 4 x_3 x_4
			- 5x_2^2
		}
		.
	\end{align*}
	Посмотрим на выражение
	\begin{equation}
		\label{eq:num_08:x1:square}
		x_1^2 + 2(x_2 + x_4)\,x_1
		.
	\end{equation}
	Как дополнить его до полного квадрата, то есть до выражения вида~$(x_1 + a)^2$, возможно, с прибавлением слагаемого, не зависящего от~$x_1$? Раскроем скобки в этом гипотетическом выражении:
	\[
		(x_1 + a)^2 =
		x_1^2 + 2ax_1 + a^2
		.
	\]
	Отсюда видно, что $a$ --- это половина множителя при~$x_1$. В нашем случае множитель при~$x_1$ --- это $2(x_2 + x_4)$, так что
	\[
		a = x_2 + x_4.
	\]
	Итак, у нас есть выражение~\eqref{eq:num_08:x1:square}, которое мы записали в виде
	\[
		x_1^2 + 2ax_1,
	\]
	где~$a = x_2 + x_4$. Доабавим к нему и вычтем~$a^2$:
	\[
		x_1^2 + 2ax_1 =
		x_1^2 + 2ax_1 + a^2 - a^2
		=
		(x_1 + a)^2 - a^2
		=
		(x_1 + x_2 + x_4)^2 - (x_2 + x_4)^2
		.
	\]
	Итак,
	\begin{equation}
		\label{eq:num_08:q:x'}
		q(x) =
		(x_1 + x_2 + x_4)^2 - (x_2 + x_4)^2
		\noatt{ {}
			- 4 x_3 x_4
			- 5x_2^2
		}
		.
	\end{equation}
	Положим теперь
	\[
		\Vect{
			x_1' \\ x_2' \\ x_3' \\ x_4'
		}
		=
		\Vect{
			x_1 + x_2 + x_4 \\
			x_2 + x_4 \\
			x_3 \\
			x_4
		}
		=
		\begin{pmatrix}
			1 & 1 & 0 & 1 \\
			0 & 1 & 0 & 1 \\
			0 & 0 & 1 & 0 \\
			0 & 0 & 0 & 1
		\end{pmatrix}
		\Vect{x_1 \\ x_2 \\ x_3 \\ x_4}
		.
	\]
	Выразим~$(x_\nu)$ через~$(x_\nu')$, найдя обратную матрицу:
	\begin{multline*}
		\left( 
			\begin{array}{rrrr|rrrr}
				1 & 1 & 0 & 1  &  1 & 0 & 0 & 0 \\
				0 & 1 & 0 & 1  &  0 & 1 & 0 & 0 \\
				0 & 0 & 1 & 0  &  0 & 0 & 1 & 0 \\
				0 & 0 & 0 & 1  &  0 & 0 & 0 & 1
			\end{array}
		\right)
		\sim
		\left( 
			\begin{array}{rrrr|rrrr}
				1 & 0 & 0 & 0  &  1 &-1 & 0 & 0 \\
				0 & 1 & 0 & 1  &  0 & 1 & 0 & 0 \\
				0 & 0 & 1 & 0  &  0 & 0 & 1 & 0 \\
				0 & 0 & 0 & 1  &  0 & 0 & 0 & 1
			\end{array}
		\right)
		\nline{\sim}
		\left( 
			\begin{array}{rrrr|rrrr}
				1 & 0 & 0 & 0  &  1 &-1 & 0 & 0 \\
				0 & 1 & 0 & 0  &  0 & 1 & 0 &-1 \\
				0 & 0 & 1 & 0  &  0 & 0 & 1 & 0 \\
				0 & 0 & 0 & 1  &  0 & 0 & 0 & 1
			\end{array}
		\right)
		.
	\end{multline*}
	Итак,
	\[
		\Vect{x_1 \\ x_2 \\ x_3 \\ x_4}
		=
		\left( 
			\begin{array}{rrrr}
				1 &-1 & 0 & 0 \\
				0 & 1 & 0 &-1 \\
				0 & 0 & 1 & 0 \\
				0 & 0 & 0 & 1
			\end{array}
		\right)
		\Vect{ x_1' \\ x_2' \\ x_3' \\ x_4' }
		.
	\]
	Подставим выражения~$x_\nu$ через~$x_\nu'$ в~$q(x)$ в выражении~\eqref{eq:num_08:q:x'}:
	\begin{align*}
		q(x) &=
		(x_1')^2 - (x_2')^2 -
		4x_3'x_4' - 5(x_2' - x_4')^2
		\nlinea{=}
		(x_1')^2 - (x_2')^2 -
		4x_3'x_4'
		- 
		5
		\bigl(
			(x_2')^2 - 2x_2'x_4' + (x_4')^2
		\bigr)
		\nlinea{=}
		(x_1')^2
		-
		6(x_2')^2
		- 
		5(x_4')^2
		+
		10x_2'x_4'
		-
		4x_3'x_4'
		.
	\end{align*}




	\textstars

	Другой способ состоит в том, чтобы посмотреть на выражение
	\[
		(a + b + c)^2 =
		a^2 + b^2 + c^2
		+
		2(ab + ac + bc)
		.
	\]
	
	\textstars

	Посмотрим на слагаемые, содержащие~$x_1$. Среди них есть два с перекрёстными слагаемыми: $2x_1 x_2$ и $2x_1 x_4$. 

	Для каждой пары перекрёстных слагаемых~$x_i x_j$ произведём преобразование
	\[
		\Vect{
			x_i \\ x_j
		}
		\mapsto
		\Vect{
			x_i' \\ x_j'
		}
		:=
		\frac{1}{2}
		\Vect{
			x_i + x_j
			\\
			x_i - x_j
		}
		,
	\]
	равносильное повороту на~$\pi/4$ в плоскости~$\pi_{ij} := \opspan(e_i, e_j)$. При этом для $\mu \notin\set*{i, j}$ имеем~$x_\mu' = x_\mu$ (остальные координаты не меняются).

\end{proofm}
